\section{Introduction}
\label{sec:introduction}

During earthquake loading, saturated granular soils can experience a rapid rise in pore pressure that drastically reduces their shear strength. This phenomenon, termed liquefaction, has caused severe damage to infrastructure and buildings \citep{IshiharaKoga1981,SeedIdriss1967}. A large body of laboratory work has investigated liquefaction resistance using cyclic undrained triaxial tests, examining the effects of cyclic stress ratio (CSR), relative density, and confining pressure \citep{Hyodo1991,SeedLee1966,Silver1976,Toki1986}. However, vertically propagating shear waves in the ground induce gradually varying shear stress on soil elements, leading to continuous rotation of principal stress axes \citep{Arthur1980,IshiharaTowhata1983}, which cannot be fully represented in conventional triaxial loading.

Hollow cylinder apparatus (HCA) and hollow torsional shear tests address this limitation by allowing torsional loading with principal stress axis rotation. Early studies using HCA-type devices \citep[e.g.,][]{IshiharaYasuda1975,Tatsuoka1986,YamashitaToki1993} showed that liquefaction behavior can differ from triaxial-test results and depends on the loading mode and stress path. These observations motivate the use of HCA-type boundary conditions when investigating liquefaction mechanisms under more realistic cyclic loading paths.

The influence of the initial lateral-to-vertical effective stress ratio, \Kzero{}, on liquefaction resistance has been investigated extensively, yet the reported findings vary considerably depending on test conditions. Some experimental studies reported weak or negligible dependence of liquefaction strength on the initial anisotropic stress state \citep[e.g.,][]{IshiharaTakatsu1979,YamashitaToki1993}, whereas others found clear differences between isotropically consolidated (IC) and anisotropically consolidated (AC) specimens, with the trend depending on density \citep{GeorgiannouKonstadinou2014}. \citet{Vargas2020} observed that specimens under compressive consolidation exhibited higher liquefaction resistance than those consolidated isotropically. More recently, \citet{Zhao2024} conducted systematic torsional shear tests covering both compressional and extensional consolidation states and attributed the higher liquefaction resistance under compressive consolidation to smaller principal stress axis rotation during cyclic loading. While these findings establish the macroscopic trend convincingly, the particle-scale mechanisms through which stress anisotropy alters liquefaction resistance remain to be explored.

The discrete element method \citep[DEM;][]{Cundall1979} is well suited for this problem because it provides direct access to particle-scale quantities such as contact topology, force chains, and fabric anisotropy, while avoiding some specimen-to-specimen variability associated with laboratory preparation. Numerous DEM studies have examined liquefaction-related behavior in triaxial, simple-shear, or polyhedral specimen settings \citep[e.g.,][]{Huang2018,Yang2021,Jiang2021,Morimoto2021,WuYang2022,Xie2023,YangHuang2023,Zhang2023} and have identified important microscopic factors governing liquefaction resistance. Recent DEM investigations have also addressed \Kzero{} effects: \citet{YangTaiebat2024} found that both preparation protocols and \Kzero{} influence liquefaction resistance, with the difference between IC and AC states narrowing at higher densities; \citet{Otsubo2022} showed that specimens with stronger inherent anisotropy exhibit weaker stiffness in the minor principal direction and correspondingly lower liquefaction resistance. However, triaxial configurations do not incorporate principal stress axis rotation, while periodic-boundary and polyhedral specimen settings, though capable of applying general stress paths, lack direct counterparts in standard laboratory apparatus. Although DEM replication of HCA tests has precedents under drained loading \citep[e.g.,][]{Li2014,Liu2021}, extension to undrained cyclic loading in HCA geometry remains limited. \citet{Han2024jgs} first realized undrained cyclic torsional shear in a DEM-based HCA model by introducing a combined servo mechanism that simultaneously satisfies the HCA stress boundary conditions and the constant-volume constraint. \citet{Ma2024} subsequently generalized this framework to accommodate a broader range of loading conditions, including drained stress paths and combined axial--torsional loading. While this framework successfully reproduced the target stress paths in most conditions, three aspects of the servo formulation can be refined.
(a)~The servo targets are formulated as absolute effective stresses on individual walls; however, these target values must in practice be inferred from the pressure difference applied between the inner and outer chambers---for example, when the chamber pressures are equal, the target effective stresses on the two cylindrical walls equal their current average, and the servo converges toward this average. The servo is therefore implicitly controlling the \emph{difference} between the inner and outer effective stresses.
(b)~Independent servo equations are written for each wall velocity ($dr/dt$, $dR/dt$, $dH/dt$), with the constant-volume constraint ($dV/dt=0$) appended as a separate condition; under undrained conditions, however, this constraint should be substituted to eliminate one kinematic variable, reducing the system to two independent variables (e.g., $dr/dt$ and $dH/dt$).
(c)~Once the volume constraint is substituted, the eliminated variable (e.g., $dR/dt$) becomes a function of the remaining two, and its contribution enters the radial and axial servo equations explicitly. Consequently, each remaining variable affects both pressure differences simultaneously, and the servo coefficients can no longer be derived from the contact stiffness at each wall independently but must be obtained by solving the coupled system.
The present study addresses these three points by (a)~adopting effective pressure differences---the radial pressure difference between the outer and inner chambers, and the axial-to-radial pressure difference controlled by the loading piston---directly as servo targets, (b)~substituting the volume constraint to eliminate one kinematic variable, and (c)~analytically inverting the resulting $2\times 2$ coupled stiffness matrix to obtain exact servo coefficients at every timestep.

Analytical exactness of the servo coefficients is important because earlier independent-gain formulations are inherently approximate: they do not account for the cross-coupling between boundary movements, so the resulting coefficients are not mathematically exact. The closed-form coupled solution derived here eliminates this approximation, ensuring exact servo coefficients at every timestep regardless of the loading stage or stress state. The full coupled formulation operates in stress-control mode during monotonic shear calibration, where both radial and axial conditions must be satisfied simultaneously; the successful reproduction of the experimental stress path and stress--strain response (Section~\ref{subsec:calibration_validation}) provides direct validation of the servo mechanism. The formulation is then applied to systematically investigate \Kzero{} effects under undrained cyclic torsional shear with direct correspondence to laboratory HCA tests.

The present work investigates the effect of initial stress anisotropy on liquefaction resistance in DEM-simulated undrained cyclic torsional shear of hollow cylindrical specimens. The investigation is organized in two stages. First, the macroscopic \Kzero{}-dependent liquefaction resistance trend is established for representative compression, isotropic, and extension states across multiple cyclic stress ratios, accompanied by energy, stiffness, and microscopic fabric analyses. Second, the simulation matrix is expanded to two relative densities and five \Kzero{} values to separate the respective roles of relative density and stress anisotropy at the particle scale.

The remainder of this paper is organized as follows. Section~\ref{sec:methodology} presents the DEM-HCA model, specimen preparation and loading protocol, calibration/validation strategy, and microscopic descriptors. Section~\ref{sec:results} reports the macroscopic and microscopic responses and discusses the mechanisms governing \Kzero{}-dependent liquefaction resistance. Section~\ref{sec:conclusions} summarizes the main findings and implications.
